% Summary of the robot commands (Python library LibRobot: from LibRobot import *)
%
% \apxCmd{...} writes the command in the first column and remembers its
% width; \apxDef, at the start of the second column, shifts the text when
% the command overflows the first column.
\newsavebox{\apxBoite}
\newcommand{\apxCmd}[1]{\global\setbox\apxBoite\hbox{\idt #1}\usebox{\apxBoite}}
\newcommand{\apxDef}{\apxFontSize\ifdim\wd\apxBoite>1.2cm
  \hspace*{\dimexpr\wd\apxBoite-1.2cm\relax}\fi}
% keyboard keys
\newcommand{\apxTouche}[1]{\mbox{\begin{picture}(12,12)(0, 2)
\put(0,0){\framebox(12,12)[lb]{}}
\put(1.5,2){\makebox(0,0)[lb]{\tiny #1}}
\end{picture}\hspace{0.3ex}}}
\newcommand{\apxEsc}{\mbox{\begin{picture}(14,12)(0, 2)
\put(0,0){\framebox(14,12)[lb]{}}
\put(1.5,2){\makebox(0,0)[lb]{\tiny Esc}}
\end{picture}\hspace{0.3ex}}}
\newcommand{\apxArret}{\apxEsc, \apxTouche{\S} or \apxTouche{\#}}
\newcommand{\apxEntree}{\framebox[11.mm][l]{$\hookleftarrow$~~~}}
\newcommand{\apxEffacement}{\framebox[11.mm][l]{$\leftarrow$~~~}}
% Python strings and characters
\newcommand{\apxChaine}[1]{{\cha \char34{}#1\char34{}}}
\newcommand{\apxCar}[1]{{\cha \textquotesingle{}#1\textquotesingle{}}}
\newcommand{\apxRobot}{{\normalsize \rm\robot}}

\noindent{\apxFontSize The commands of the \apxRobot\ are
those of the library {\idt LibRobot} (a program starts with
{\mcl from}\ {\idt LibRobot}\ {\mcl import}\ {\idt *}). A command is written followed by
parentheses, empty when it has no parameter, like {\idt fd()}; the
parameters are written between the parentheses, separated by commas, like
{\idt pw(2)}, {\idt fdl(-3.5)} or {\idt moveTo(2, 3)}. A real number (of type
{\idt float}) is written with a point, like {\idt 20.5}, an integer (of type
{\idt int}) without a point; an integer can be given where a real is expected,
like {\idt gs(20)}. A character string is of type {\idt str}, a boolean
of type {\idt bool} ({\mcl True} or {\mcl False}).
The functions {\idt val}\ldots\ return the current value of a setting of the
\apxRobot.\par}
\vspace{3mm}

\noindent\begin{tabular}{@{}|p{1cm}p{11.9cm}|@{}} \hline
\multicolumn{2}{|l|}{Grid of the stopping points of the {\normalsize \robot}:}  \\ \hline \hline
\apxCmd{dg()} & \apxDef (\underline{d}raw \underline{g}rid)
draws the grid of the stopping points of the \apxRobot\ (the
drawing of the grid can be interrupted by \apxArret). \\
\apxCmd{gs($n$)} & \apxDef (\underline{g}rid \underline{s}ize) sets the new length $n$ (a real
$n \geq 0$) of the sides of a square of the grid ($12 \times 12$ pixels
by default). \\
\apxCmd{gsX($n$)} & \apxDef sets the new length $n$ (a real) of the
horizontal side (\underline{X}) of a square of the grid (\underline{g}rid \underline{s}ize, 12 pixels by default). \\
\apxCmd{gsY($n$)} & \apxDef sets the new length $n$ (a real) of the
vertical side (\underline{Y}) of a square of the grid (\underline{g}rid \underline{s}ize, 12 pixels by default). \\
\apxCmd{gsRT($r$, $t$)} & \apxDef sets the new lengths of the
sides of a square of the \underline{g}rid (\underline{g}rid \underline{s}ize) by the polar coordinates \underline{$r$} and
\underline{$t$} (two reals: the length $r \geq 0$ of the diagonal of the square is
expressed in pixels and the angle $t$ between this diagonal and the horizontal side is
expressed in degrees, usually $0 \leq t \leq 90$).
\\ \hline
\apxCmd{valGsX()} & \apxDef is equal to the \underline{val}ue (of type {\idt float}) of the
length (\underline{g}rid \underline{s}ize) of the horizontal side (\underline{X}) of a
square of the grid (counted in pixels). \\
\apxCmd{valGsY()} & \apxDef is equal to the \underline{val}ue (of type {\idt float}) of the
length (\underline{g}rid \underline{s}ize) of the vertical side (\underline{Y}) of a
square of the grid (counted in pixels). \\
\apxCmd{valGsR()} & \apxDef is equal to the \underline{val}ue (of type {\idt float}) of the
length (\underline{g}rid \underline{s}ize) of the diagonal (\underline{R}) of a square
of the grid (counted in pixels). \\
\apxCmd{valGsT()} & \apxDef is equal to the \underline{val}ue (of type {\idt float}) of
the angle (\underline{T}) of the diagonal of a square
of the \underline{g}rid with its horizontal side (counted in degrees). \\
\hline \end{tabular}

\noindent\begin{tabular}{@{}|p{1cm}p{11.9cm}|@{}} \hline
\apxCmd{valX()} & \apxDef is equal to the \underline{val}ue (of type {\idt float}) of
the absolute abscissa (\underline{X}) of the \apxRobot\ (counted in pixels from the left
edge of the window). \\
\apxCmd{valY()} & \apxDef is equal to the \underline{val}ue (of type {\idt float}) of
the absolute ordinate (\underline{Y}) of the \apxRobot\ (counted in pixels from the bottom of
the window), so that after {\idt x = round(valX())} and {\idt y = round(valY())},
later on {\idt pu();\ drawTo(x, y)} puts the \apxRobot\ back at this place on
the screen. \\ \hline
\end{tabular}

\noindent\begin{tabular}{@{}|p{1cm}p{11.9cm}|@{}} \hline
\multicolumn{2}{|l|}{Relative turns on the spot:}  \\ \hline \hline
\apxCmd{etr()} & \apxDef (\underline{e}ighth \underline{t}urn
\underline{r}ight) turns the \apxRobot\ on the spot an eighth of a turn
to the right. \\
\apxCmd{qtr()} & \apxDef (\underline{q}uarter \underline{t}urn
\underline{r}ight) turns the \apxRobot\ on the spot a quarter turn
to the right. \\
\apxCmd{tetr()} & \apxDef (\underline{t}hree \underline{e}ighths of a
\underline{t}urn \underline{r}ight) turns the \apxRobot\ on the
spot 3 eighths of a turn to the right. \\
\apxCmd{etl()} & \apxDef (\underline{e}ighth \underline{t}urn
\underline{l}eft) turns the \apxRobot\ on the spot an eighth of a turn
to the left. \\
\apxCmd{qtl()} & \apxDef (\underline{q}uarter \underline{t}urn
\underline{l}eft) turns the \apxRobot\ on the spot a quarter turn
to the left. \\
\apxCmd{tetl()} & \apxDef (\underline{t}hree \underline{e}ighths of a
\underline{t}urn \underline{l}eft) turns the \apxRobot\ on the
spot 3 eighths of a turn to the left. \\
\apxCmd{ht()} & \apxDef (\underline{h}alf \underline{t}urn)
turns the \apxRobot\ on the spot a half turn. \\ \hline
\end{tabular}

\noindent\begin{tabular}{@{}|p{1cm}p{11.9cm}|@{}} \hline
\multicolumn{2}{|l|}{Absolute turns on the spot: }  \\ \hline \hline
\apxCmd{fn()} & \apxDef (\underline{f}ace \underline{n}orth) turns the \apxRobot\ on the spot towards the north (\underline{N}). \\
\apxCmd{fne()} & \apxDef makes the \apxRobot\ \underline{f}ace the north-east (\underline{NE}), turning on the spot. \\
\apxCmd{fe()} & \apxDef makes the \apxRobot\ \underline{f}ace the east (\underline{E}), turning on the spot. \\
\apxCmd{fse()} & \apxDef makes the \apxRobot\ \underline{f}ace the south-east (\underline{SE}), turning on the spot. \\
\apxCmd{fs()} & \apxDef makes the \apxRobot\ \underline{f}ace the south (\underline{S}), turning on the spot. \\
\apxCmd{fsw()} & \apxDef makes the \apxRobot\ \underline{f}ace the south-west (\underline{SW}), turning on the spot. \\
\apxCmd{fw()} & \apxDef makes the \apxRobot\ \underline{f}ace the west (\underline{W}), turning on the spot. \\
\apxCmd{fnw()} & \apxDef makes the \apxRobot\ \underline{f}ace the north-west (\underline{NW}), turning on the spot. \\
\apxCmd{face($o$)} & \apxDef turns the \apxRobot\ on the spot to \underline{face}
the \underline{o}rientation given by the value of the integer $o$ (between 0 and 7) equal to
{\idt north} $=$ 0, {\idt northEast} $=$ 1, {\idt east} $=$ 2, {\idt southEast} $=$ 3,
{\idt south} $=$ 4, {\idt southWest} $=$ 5, {\idt west} $=$ 6 or {\idt northWest} $=$ 7. \\
\apxCmd{rot($t$)} & \apxDef changes the coordinate system and the orientation of the \apxRobot\ so
that it performs a \underline{rot}ation on the spot by an angle of $t$ degrees ($t$ real,
counterclockwise if $t > 0$ and clockwise if
$t < 0$, like {\idt rot(15)} or {\idt rot(-30)}). \\ \hline
\apxCmd{valOrientation()} & \apxDef is equal to the \underline{val}ue (of type {\idt int})
of the \underline{orientation} of the \apxRobot\ coded by {\idt north} $=$ 0, {\idt northEast} $=$ 1,
{\idt east} $=$ 2, {\idt southEast} $=$ 3, {\idt south} $=$ 4, {\idt southWest} $=$ 5,
{\idt west} $=$ 6 and {\idt northWest} $=$ 7. The directions {\idt north}, {\idt east},
{\idt south} and {\idt west} are parallel to the edges of the screen, {\idt north} at the top.
The directions {\idt northEast}, {\idt southEast}, {\idt southWest} and {\idt northWest}
correspond to the orientation of the diagonals of a square of the grid.\\ \hline
\end{tabular}

\noindent\begin{tabular}{@{}|p{1cm}p{11.9cm}|@{}} \hline
\multicolumn{2}{|l|}{Pen: }  \\ \hline \hline
\apxCmd{pd()} & \apxDef (\underline{p}en \underline{d}own) puts the pen down to draw while moving. \\
\apxCmd{pu()} & \apxDef (\underline{p}en \underline{u}p) lifts the pen to avoid drawing while moving. \\
\apxCmd{dc()} & \apxDef (\underline{d}raw a \underline{c}ross) draws a small cross (as thick as the
pen), without moving the \apxRobot. \\
\apxCmd{dp()} & \apxDef (\underline{d}raw a \underline{p}oint) draws a dot (as thick as the
pen), without moving the \apxRobot. \\
\apxCmd{pw($n$)} & \apxDef (\underline{p}en \underline{w}idth) sets the new thickness $n$ (an integer) of the
pen (1 pixel by default). \\ \hline
\apxCmd{valPw()} & \apxDef is equal to the \underline{val}ue (of type {\idt int}) of
the thickness of the pen (\underline{p}en \underline{w}idth). \\
\apxCmd{valPd()} & \apxDef is equal to the boolean \underline{val}ue {\mcl True} if
and only if the pen is down (\underline{p}en \underline{d}own), otherwise to {\mcl False}. \\ \hline
\end{tabular}

\noindent\begin{tabular}{@{}|p{1cm}p{11.9cm}|@{}} \hline
\multicolumn{2}{|l|}{Moves (drawing or not depending on whether the pen is up or down): }  \\ \hline \hline
\apxCmd{fd()} & \apxDef (\underline{f}orwar\underline{d}) moves the \apxRobot\ straight
ahead to the next stopping point on the grid. \\
\apxCmd{cr()} & \apxDef (\underline{c}urve \underline{r}ight) makes the \apxRobot\ turn
by a bend of a quarter circle to the right. \\
\apxCmd{cl()} & \apxDef (\underline{c}urve \underline{l}eft) makes the \apxRobot\ turn
by a bend of a quarter circle to the left. \\
\apxCmd{fdn($n$)} & \apxDef (\underline{f}orwar\underline{d} \underline{n} times) moves the
\apxRobot\ straight ahead $n$ times ($n$ an integer; doing nothing if $n = 0$ and moving
backwards if $n < 0$, like {\idt fdn(-2)}). \\
\apxCmd{fdl($l$)} & \apxDef (\underline{f}orwar\underline{d} by a \underline{l}ength) changes the coordinate system (but not the angle of the directions NE-E)
and moves the \apxRobot\ forward by a translation of $l$ pixels ($l$ real)
in the direction it is facing (if $l > 0$) or in the opposite direction (if $l < 0$).
\\
\apxCmd{drawToReal($x$, $y$)} & \apxDef (\underline{draw to} the point with \underline{real}
coordinates) draws a straight line
segment (if the pen is down) from the current position of the \apxRobot\ to the point with
real coordinates $(x, y)$ of the Cartesian coordinate system of the grid (see
{\idt moveTo}) and moves the \apxRobot\ to this point. \\
\apxCmd{drawTo($x$, $y$)} & \apxDef does the same up to the point with integer coordinates
$(x, y)$ counted in pixels from the lower left corner of the window (like
{\idt valX()} and {\idt valY()}).\\ \hline
\apxCmd{valFdl()} & \apxDef is equal to the \underline{val}ue (of type {\idt float}) of the
length of the segment drawn by the last move forward by {\idt fd()} or by
{\idt fdl()} (\underline{f}orwar\underline{d} by a \underline{l}ength), counted in pixels. \\ \hline
\end{tabular}

\noindent\begin{tabular}{@{}|p{1cm}p{11.9cm}|@{}} \hline
\multicolumn{2}{|l|}{Positioning at the centre and on the edges of the frame: }  \\ \hline \hline
\apxCmd{gc()} & \apxDef (\underline{g}o to the \underline{c}entre) places the \apxRobot\ at the centre of {\em the screen}, which
is also {\em the centre of the grid}, without changing the orientation of the \apxRobot\ and
without drawing anything. \\
\apxCmd{gt()} & \apxDef (\underline{g}o to the \underline{t}op) places the
\apxRobot\ at the top of the frame, without changing its horizontal position or its
orientation and without drawing anything. \\
\apxCmd{gb()} & \apxDef (\underline{g}o to the \underline{b}ottom) places the
\apxRobot\ at the bottom of the frame, without changing its horizontal position or its
orientation and without drawing anything. \\
\apxCmd{gl()} & \apxDef (\underline{g}o to the \underline{l}eft) places the
\apxRobot\ on the left of the frame, without changing its vertical position or its
orientation and without drawing anything. \\
\apxCmd{gr()} & \apxDef (\underline{g}o to the \underline{r}ight) places the
\apxRobot\ on the right of the frame, without changing its vertical position or its
orientation and without drawing anything. \\ \hline
\end{tabular}

\noindent\begin{tabular}{@{}|p{1cm}p{11.9cm}|@{}} \hline
\multicolumn{2}{|l|}{Cartesian coordinate system: }  \\ \hline \hline
\apxCmd{translate($x$, $y$)} & \apxDef \underline{translate}s the \apxRobot\ by $x$
squares of the grid horizontally (to the right if $x > 0$ and to the left if $x < 0$) and $y$
squares of the grid vertically (up if $y > 0$ and down if $y < 0$), without changing
its orientation and without drawing anything ($x$ and $y$ integers). \\
\apxCmd{moveTo($x$, $y$)} & \apxDef (\underline{move to}) moves the
\apxRobot\ to the point with integer coordinates $(x, y)$ on the grid, without changing its
orientation or drawing anything, the origin of the Cartesian coordinate system being at the
centre of the grid. \\
\apxCmd{moveToReal($x$, $y$)} & \apxDef does the same with \underline{real} coordinates
$x$ and $y$ (of type {\idt float}), like {\idt moveToReal(1.5, -2.5)}. \\ \hline
\end{tabular}

\noindent\begin{tabular}{@{}|p{1cm}p{11.9cm}|@{}} \hline
\multicolumn{2}{|l|}{Boolean test of the orientation of the {\normalsize \rm\robot}:  }  \\ \hline \hline
\apxCmd{tn()} & \apxDef (\underline{t}est \underline{n}orth) \underline{t}ests whether the \apxRobot\ is facing north (\underline{N}). \\
\apxCmd{tne()} & \apxDef \underline{t}ests whether the \apxRobot\ is facing north-east (\underline{NE}). \\
\apxCmd{te()} & \apxDef \underline{t}ests whether the \apxRobot\ is facing east (\underline{E}).\\
\apxCmd{tse()} & \apxDef \underline{t}ests whether the \apxRobot\ is facing south-east (\underline{SE}).\\
\apxCmd{ts()} & \apxDef \underline{t}ests whether the \apxRobot\ is facing south (\underline{S}).\\
\apxCmd{tsw()} & \apxDef \underline{t}ests whether the \apxRobot\ is facing south-west (\underline{SW}).\\
\apxCmd{tw()} & \apxDef \underline{t}ests whether the \apxRobot\ is facing west (\underline{W}).\\
\apxCmd{tnw()} & \apxDef \underline{t}ests whether the \apxRobot\ is facing north-west (\underline{NW}). \\ \hline
\end{tabular}

\noindent\begin{tabular}{@{}|p{1cm}p{11.9cm}|@{}} \hline
\multicolumn{2}{|l|}{Boolean test of the position of the {\normalsize \rm\robot}:  }  \\ \hline \hline
\apxCmd{tb()} & \apxDef (\underline{t}est \underline{b}order) \underline{t}ests whether
the \apxRobot\ is at the edge (\underline{b}order) of the frame (a move of the \apxRobot\
forward by {\idt fd()} would take it out of the frame). \\ \hline
\end{tabular}

\noindent\begin{tabular}{@{}|p{1cm}p{11.9cm}|@{}} \hline
\multicolumn{2}{|l|}{Time:  }  \\ \hline \hline
\apxCmd{clock()} & \apxDef returns three integers: the present hour (between 0
and 23), minute (between 0 and 59) and second (between 0 and 59), so that after
{\idt h, mi, s = clock()} the variables {\idt h}, {\idt mi} and {\idt s} contain
these three values. \\
\apxCmd{date()} & \apxDef likewise returns three integers, the day (between 1 and 31), the month
(between 1 and 12) and the current year, as in {\idt d, mo, y = date()}. \\ \hline
\apxCmd{delay($t$)} & \apxDef waits for a \underline{delay} of $t$ seconds ($t$ of type
{\idt float}, like {\idt delay(0.5)}). \\
\apxCmd{tickCount()} & \apxDef is equal to the integer number of sixtieths of a second
elapsed (to measure a duration as the difference of two values). \\ \hline
\end{tabular}

\noindent\begin{tabular}{@{}|p{1cm}p{11.9cm}|@{}} \hline
\multicolumn{2}{|l|}{Interrupting the execution of the program:  }  \\ \hline \hline
\apxCmd{\rm\apxArret} & \apxDef typed while the \apxRobot\ is moving (by
{\idt fd()}, {\idt cr()}, {\idt cl()}, \ldots) or running step by step, any of these
keys stops the program and closes the window. \\ \hline
\end{tabular}

\noindent\begin{tabular}{@{}|p{1cm}p{11.9cm}|@{}} \hline
\multicolumn{2}{|l|}{ Debugging the program:  }  \\ \hline \hline
\apxCmd{st()} & \apxDef \underline{st}ops the \apxRobot\ and recalls in the bottom line
of the window the keys that can be used: \apxTouche{$+$} increases and \apxTouche{$-$}
decreases the speed of the \apxRobot\ (between 0 and 10, speed 0 switching to step-by-step
mode), \apxEsc\ or \apxTouche{\#} closes the window and ends the program, any other
key redraws the whole drawing from the beginning, at the chosen speed. A program
that ends without {\idt st()} waits for a key to be typed or a click before
closing the window. \\
\apxCmd{ss()} & \apxDef (\underline{s}tep by \underline{s}tep) runs the
rest of the program in step-by-step mode, at speed 0: the \apxRobot\ waits for a key to be
typed before each command (\apxArret\ stopping the program). \\
\apxCmd{ex()} & \apxDef runs the rest of the program in normal \underline{ex}ecution mode,
at speed 10. \\
\apxCmd{sp($v$)} & \apxDef runs the rest of the program at \underline{sp}eed
$v$ (an integer; step by step if $v \leq 0$, with a maximum of 10 if $v > 10$). \\
\apxCmd{showRobot()} & \apxDef makes the \apxRobot\ visible on the screen
during the drawing (default option). \\
\apxCmd{hideRobot()} & \apxDef makes the \apxRobot\ invisible on the screen
during the drawing. \\
\apxCmd{forbidExits()} & \apxDef forbids the \apxRobot\ to go out
of the frame in the rest of the program (default option: a command
that would take the \apxRobot\ out of the frame then stops the program with an error,
writing in the terminal, for example, {\tt Error line 4 of Program.py: the robot
would leave the frame (fd)}; the window stays displayed). \\
\apxCmd{allowExits()} & \apxDef allows the \apxRobot\ to go out
of the frame in the rest of the execution of the program. \\ \hline
\apxCmd{valSp()} & \apxDef is equal to the \underline{val}ue (of type {\idt int}) of the
\underline{sp}eed of the \apxRobot\ (between 0 (step-by-step mode)
and 10). \\
\apxCmd{valVisible()} & \apxDef is equal to the boolean value {\mcl True} (respectively
{\mcl False}) if the \apxRobot\ is \underline{visible} (respectively invisible) on the screen. \\
\apxCmd{valAllowExits()} & \apxDef is equal to the boolean value {\mcl True}
(respectively {\mcl False}) if exits out of the frame are
allowed (\underline{allow exits}) (respectively forbidden). \\ \hline
\end{tabular}

\noindent\begin{tabular}{@{}|p{1cm}p{11.9cm}|@{}} \hline
%
Colour: & \\ \hline \hline
%
\multicolumn{2}{|p{12.9cm}|}{Colours are integers: {\idt white} $= 0$,
{\idt yellow} $= 1$, {\idt blue} $= 2$, {\idt violet} $= 3$,
{\idt green} $= 4$, {\idt red} $= 5$, {\idt indigo} $= 6$ and
{\idt black} $= 7$.}\\\hline \hline
%
\apxCmd{pc($c$)} & \apxDef (\underline{p}en \underline{c}olour) sets the colour $c$ of the pen, like {\idt pc(red)}. \\
\apxCmd{paint()} & \apxDef paints with the colour of the pen the region
of a single colour (usually white), bounded by the lines of the drawing, where the
\apxRobot\ is. \\
\apxCmd{bg($c$)} & \apxDef (\underline{b}ack\underline{g}round colour) completely erases the drawing then fills the screen with the background colour $c$.
The \apxRobot\ is placed at the centre of the screen
in the initial conditions (visible, facing north, exits out of the frame
forbidden, black pen of thickness 1 down and normal execution at speed 10); the
dimensions of the grid do not change. \\ \hline
\apxCmd{valBg()} & \apxDef is equal to the \underline{val}ue (of type {\idt int}) of the
background colour (\underline{b}ack\underline{g}round). \\
\apxCmd{valPc()} & \apxDef is equal to the \underline{val}ue (of type {\idt int}) of the
pen colour (\underline{p}en \underline{c}olour). \\
\apxCmd{blackAndWhiteScreen()} & \apxDef is equal to {\mcl True} if
and only if the screen is black and white, otherwise to {\mcl False} (it is always
{\mcl False} on today's screens). \\ \hline
\end{tabular}

\noindent\begin{tabular}{@{}|p{1cm}p{11.9cm}|@{}} \hline
\multicolumn{2}{|l|}{Writing in the drawing window of the {\normalsize \rm\robot}:  }  \\ \hline \hline
\apxCmd{label($c$)} & \apxDef writes the character string $c$ (between quotes
\apxChaine{} and \apxChaine{}, where any quote must be preceded by a
backslash: {\cha \textbackslash\char34}, like {\idt label(}\apxChaine{Hello}{\idt )})
to the right of the \apxRobot\ (or to its left if it is
facing west), without changing its position or that of its pen. \\
\apxCmd{intToString($i$)} & \apxDef (\underline{int}eger \underline{to} \underline{string}) is the character string
representing the integer $i$ in base 10 (like {\idt str($i$)}). \\
\apxCmd{realToString($r$)} & \apxDef (\underline{real} \underline{to} \underline{string}) is the character string
representing the real $r$ in the form {\idt $d.dddd$e$\pm ee$} where $d$ is a digit of the
mantissa, $e$ is a digit of the exponent and {\idt e} reads ``times 10 to the power''. \\
\apxCmd{boolToString($b$)} & \apxDef (\underline{bool}ean \underline{to} \underline{string}) is the character string
(\apxChaine{True} or \apxChaine{False}) re\-pre\-sent\-ing the boolean $b$.\\
\apxCmd{concat([$c_1$, $c_2$, \ldots, $c_n$])} & \apxDef is the character
string obtained by \underline{concat}enating the strings $c_1$, $c_2$, \ldots, $c_n$ one
after the other (like $c_1$ {\idt +} $c_2$ {\idt +} \ldots\
{\idt +} $c_n$). \\ \hline
\apxCmd{textFont($f$)} & \apxDef chooses the \underline{font}
$f$ of the \underline{text} ({\idt courier}, {\idt times}, {\idt helvetica}, \ldots). \\
\apxCmd{textSize($s$)} & \apxDef chooses the \underline{size}
$s$ (an integer) of the \underline{text}. \\
\apxCmd{textStyle($s$)} & \apxDef chooses the \underline{style} of the
\underline{text} ({\idt normal}, {\idt bold}, {\idt italic},
{\idt underline}, \ldots). \\
\multicolumn{2}{|p{12.9cm}|}{\apxFontSize These three commands have no visible effect: the
font of the window of the \apxRobot\ cannot be changed.} \\ \hline
\apxCmd{valTextFont()} & \apxDef \underline{val}ue (of type {\idt str}) of the
current \underline{font} of the \underline{text}. \\
\apxCmd{valTextSize()} & \apxDef \underline{val}ue (of type {\idt int}) of the
current \underline{size} of the \underline{text}. \\ \hline
\end{tabular}

\noindent\begin{tabular}{@{}|p{1cm}p{11.9cm}|@{}} \hline
\multicolumn{2}{|l|}{Messages displayed in the bottom line of the window of the {\normalsize \rm\robot}:  }  \\ \hline \hline
\apxCmd{message($c$)} & \apxDef erases the \underline{message} in the bottom line of the
window of the \apxRobot\ then writes the character string $c$ there. \\
\apxCmd{markPause()} & \apxDef (\underline{mark} a \underline{pause}) pauses (making the symbol
$\hookleftarrow$ blink in the message line) until the mouse button or a key of the
keyboard is pressed.\\
\apxCmd{clearMessage()} & \apxDef (\underline{clear} the \underline{message}) erases the message in the
bottom line of the window of the \apxRobot. \\ \hline
\apxCmd{readChar()} & \apxDef (\underline{read} a \underline{char}acter) waits (making the symbol $\hookleftarrow$ blink
in the message line) until a key of the keyboard or the mouse button is
pressed, then reads and returns the character (a string of a single
character, of type {\idt str}) corresponding to the key pressed (or the character
\apxCar{0} if the mouse button was pressed). \\ \hline
\apxCmd{keyboardCommandsMemento} & \apxDef is a character string
recalling the \underline{keyboard} \underline{commands} for interactive control of the \apxRobot\ (see {\idt interaction}), to
be displayed by {\idt message(keyboardCommandsMemento)}. \\
\apxCmd{valMessage()} & \apxDef is equal to the \underline{val}ue (of type {\idt str})
of the \underline{message} written in the bottom line of the window of the \apxRobot\ or to the
empty string \apxChaine{} if there is no message. \\ \hline
\end{tabular}
\vfil
\noindent\begin{tabular}{@{}|p{1cm}p{11.9cm}|@{}} \hline
\multicolumn{2}{|l|}{Reading and writing text:  }  \\ \hline \hline
\multicolumn{2}{|p{12.9cm}|}{\apxFontSize The text written by {\idt write} and
{\idt writeln} is displayed in the terminal. When the program waits for an answer
({\idt readInt()}, \ldots), the question (the current line, preceded by the last two
lines written) is displayed in the bottom line of the window of the \apxRobot. The
answer is typed in the window of the \apxRobot, where it is displayed after the question
(\apxEffacement\ to correct, \apxEntree\ to finish), or else in the terminal.} \\ \hline
\apxCmd{write($c_1$, \ldots, $c_n$)} & \apxDef writes the character strings (or the
numbers) $c_1$, \ldots, $c_n$ one after the other, like
{\idt write(\apxChaine{N = }, n)}. \\
\apxCmd{writeln($c_1$, \ldots, $c_n$)} & \apxDef does the same then goes to a new line
({\idt writeln()} only goes to a new line). \\ \hline
\apxCmd{readInt()} & \apxDef (\underline{read} an \underline{int}eger) reads the user's answer, ended by
\apxEntree, and returns the integer (of type {\idt int}) that it represents. If the
answer is not an integer, the question \apxChaine{Enter an integer: } is asked again.
On a French keyboard, the digits can be typed without the Shift
key. \\
\apxCmd{readReal()} & \apxDef does the same to \underline{read} a
\underline{real} number (of type {\idt float}). \\
\apxCmd{readString()} & \apxDef (\underline{read} a \underline{string}) reads the user's answer, ended by
\apxEntree, and returns it as a character string (of type {\idt str}). \\
\apxCmd{readCharacter()} & \apxDef does the same and returns the first
\underline{character} of the answer (a string of a single character). \\
\apxCmd{readln()} & \apxDef waits until the user types \apxEntree\ (what was typed
before is ignored). \\
\apxCmd{clearText()} & \apxDef (\underline{clear} the \underline{text}) erases the text of the terminal. \\ \hline
\end{tabular}
\vfil
\noindent\begin{tabular}{@{}|p{1cm}p{11.9cm}|@{}} \hline
\multicolumn{2}{|l|}{Useful functions:  }  \\ \hline \hline
\apxCmd{pi} & \apxDef is the real $\pi = 3.14159\ldots$ (of type {\idt float}). \\
\apxCmd{round($x$)} & \apxDef is the integer closest to the real $x$, like
{\idt round(2.5)} $=$ 3 and {\idt round(-2.5)} $= -3$ (the library replaces the
{\idt round} of \python, which rounds halves to the even number: 2 for 2.5). \\
\apxCmd{random()} & \apxDef is an integer drawn at random between $-32767$ and $32767$. \\ \hline
\end{tabular}
\vfil
\noindent\begin{tabular}{@{}|p{1cm}p{11.9cm}|@{}} \hline
\multicolumn{2}{|l|}{Sounds:  }  \\ \hline \hline
\apxCmd{beep()} & \apxDef emits a short sound. \\
\apxCmd{longBeep()} & \apxDef emits a long sound. \\ \hline
\end{tabular}
\vfil
\noindent\begin{tabular}{@{}|p{1cm}p{11.9cm}|@{}} \hline
\multicolumn{2}{|l|}{Music:  }  \\ \hline \hline
\apxCmd{metronome($n$)} & \apxDef sets the \underline{metronome} to $n$ quarter notes per minute
({\idt largo} $=$ 50, {\idt lento} $=$ 56, {\idt larghetto} $=$ 60, {\idt adagio} $=$ 70,
{\idt andante} $=$ 85, {\idt andantino} $=$ 95, {\idt moderato} $=$ 110,
{\idt allegretto} $=$ 115, {\idt allegro} $=$ 130, {\idt vivace} $=$ 150,
{\idt presto} $=$ 175, {\idt prestissimo} $=$ 200; 100 by default). \\
\apxCmd{dynamics($n$)} & \apxDef sets the \underline{dynamics} $n$ (between 1 and 9:
{\idt pianississimo} $=$ 1, {\idt pianissimo} $=$ 2, {\idt piano} $=$ 3,
{\idt mezza\_voce} $=$ 4, {\idt mezzo\_piano} $=$ 4, {\idt mezzo\_forte} $=$ 6,
{\idt forte} $=$ 7, {\idt fortissimo} $=$ 8, {\idt fortississimo} $=$ 9). \\
\apxCmd{note($p$, $d$)} & \apxDef plays a \underline{note} of pitch
$p = o * \mbox{\idt octave} + n + a$ (an integer) where $o = -3$, $-2$, \ldots, $3$ is the octave
of the note (0 for the octave of middle C of the piano, {\idt octave} $=$ 12), $n$ $=$
{\idt C} $=$ 0, {\idt D} $=$ 2, {\idt E} $=$ 4, {\idt F} $=$ 5,
{\idt G} $=$ 7, {\idt A} $=$ 9, {\idt B} $=$ 11 is the name of the note,
$a$ $=$ {\idt flat} $= -1$, {\idt natural} $= 0$, {\idt sharp} $= 1$ is the accidental of the note
and $d$ (a real) is the duration of the note, counted in quarter notes: {\idt whole}
$= 4.$, {\idt dotted\_whole} $= 6.$, {\idt half} $= 2.$, {\idt dotted\_half} $= 3.$,
{\idt half\_triplet} $= 4/3$, {\idt quarter} $= 1.$,
{\idt dotted\_quarter} $= 1.5$, {\idt quarter\_triplet} $= 2/3$, {\idt eighth}
$= 0.5$, {\idt dotted\_eighth} $= 0.75$,
{\idt eighth\_triplet} $= 1/3$, {\idt sixteenth} $= 0.25$,
{\idt dotted\_sixteenth} $= 0.375$ or {\idt thirty\_second}
$= 0.125$, like {\idt note(octave + G + sharp, eighth)}. \\
\apxCmd{rest($d$)} & \apxDef makes a \underline{rest} of duration $d$ $=$
{\idt whole\_rest} $= 4.$, {\idt dotted\_half\_rest} $= 3.$, {\idt half\_rest}
$= 2.$, {\idt dotted\_quarter\_rest} $= 1.5$, {\idt quarter\_rest} $= 1.$,
{\idt dotted\_eighth\_rest} $= 0.75$, {\idt eighth\_rest} $= 0.5$,
{\idt sixteenth\_rest} $= 0.25$ or
{\idt thirty\_second\_rest} $= 0.125$. \\ \hline
\end{tabular}
\vfil
\noindent\begin{tabular}{@{}|p{1cm}p{11.9cm}|@{}} \hline
\multicolumn{2}{|l|}{Interaction with the user of the {\normalsize \rm\robot}:  }  \\ \hline \hline
\multicolumn{2}{|l|}{\idt interaction(coordinates, turns, drawing, mouseEnds, save)}\\
& \apxFontSize erases the screen then lets the user move the \apxRobot\ with the mouse,
by clicking on a point of the grid (if {\idt mouseEnds} is {\mcl True}: from square to square, button released)
or with the arrow keys, and possibly drive it with the keyboard (the five
parameters are
booleans, {\mcl True} or {\mcl False}). The interaction ends with the key \apxEntree\ (or by
pressing the mouse button if {\idt mouseEnds} is {\mcl True}); the program
then stops as with {\idt st()}.

\tiret If {\idt coordinates} is {\mcl True}, the coordinates of the \apxRobot\ are displayed
during the interaction.

\tiret If {\idt turns} is {\mcl True}, the user can turn the
\apxRobot\ with the keys \apxTouche{$-$} ({\idt qtl}), \apxTouche{$+$} ({\idt qtr}),
\apxTouche{l} ({\idt etl}) and \apxTouche{r} ({\idt etr}).

\tiret If {\idt turns} and {\idt drawing} are {\mcl True}, the user can also draw with
the keys \apxTouche{f} ({\idt fd}), \apxTouche{$($} ({\idt cl}), \apxTouche{$)$}
({\idt cr}), \apxTouche{u} ({\idt pu}), \apxTouche{d} ({\idt pd}), \apxTouche{p}
({\idt dp}), \apxTouche{c} ({\idt dc}), \apxTouche{$.$} ({\idt dg}), \apxTouche{$*$}
(change of the pen colour, {\idt pc}) and \apxTouche{w} ({\idt paint}), undoing the
last command with \apxEffacement.

\tiret If {\idt drawing} and {\idt save} are {\mcl True}, the program asks at the end
whether to save the drawing: by answering \apxTouche{y}, one obtains a program
{\idt my\_drawing.py} whose execution reproduces the drawing.
\\ \hline
\multicolumn{2}{|l|}{\idt x, y, c = interactionTP(coordinates, turns, drawing,}\\
\multicolumn{2}{|l|}{\idt \phantom{x, y, c = interactionTP(}mouseEnds, save)}\\
& \apxFontSize lets the user move the \apxRobot\ with the mouse (or the arrow keys)
and drive it with the keyboard like {\idt interaction} (except \apxTouche{$*$}, \apxTouche{w} and
\apxEffacement), but without erasing the screen, and then continues the
program. The interaction ends with the key \apxEntree, with \apxEsc, with any key
that is not a command of the \apxRobot\ or, if {\idt mouseEnds} is {\mcl True}, by
pressing the mouse button. At the end, {\idt interactionTP} returns three values, put
here in the variables {\idt x}, {\idt y} and {\idt c}: the integers {\idt x} and {\idt y}
are the horizontal and vertical displacement of the \apxRobot\ counted in squares of the grid,
and {\idt c} is the last key typed (a space \apxCar{ } if
the interaction ended with a click). The parameter {\idt save} has no effect.
\\ \hline
\end{tabular}
\newpage                 % If the chapter ends with a blank page
\thispagestyle{empty}
